When a language-model agent improves after receiving a ``skill'' --- a human-written
or self-evolved procedure --- the gain can mean two very different things: the skill
injected genuinely new, transferable knowledge the model lacked, or it merely
re-elicited latent knowledge the model already had, a dressed-up prompt optimization.
Existing skill ablations report only task pass-rate and cannot tell these apart.
We build a controlled substrate in which the answer to ``novel or latent?'' is known by
construction --- counterfactual arithmetic (base-9/11) and an invented-language glossary
(``\zeth'') whose word mappings are randomly generated at run time and thus cannot appear
in any training corpus --- and we propose the Elicitation--Novelty Discriminator (\END),
built from three measurable quantities: self-generation utility (can the base model write a
working skill unaided?), length-controlled gain (does the skill beat a token-matched
irrelevant control?), and whether verifier-only self-evolution can reach the skill.
Across two frontier models and 3{,}000+ deterministic API evaluations, a skill injects
genuinely new knowledge in exactly one regime: where the model provably cannot self-generate
it. On \zeth the correct skill lifts accuracy $0.00 \to 1.00$ (both models, $p<10^{-4}$,
Cohen's $h=3.14$) while the self-generated skill and a length-matched control stay at $0.00$;
a five-step self-evolution loop never moves \zeth off $0.00$ under verifier-only feedback but
reaches $0.85$ the instant it receives ground-truth answers. Self-generated skills can even
harm: on base-9 multiplication a model's own skill dropped accuracy $0.85 \to 0.25$. We conclude
that current self-evolving skills are, rigorously, better elicitation bounded by the base
model's latent knowledge --- not autonomous discovery --- and we deliver a reusable protocol
for telling the two apart.
